\mysection{assert()の呼び出し}
\myindex{\CStandardLibrary!assert()}

時々、\TT{assert()}マクロの存在も有用です：
一般的に、このマクロはソースファイル名、行番号、条件をコードに残します。

最も有用な情報はassertの条件に含まれており、そこから変数名や構造体フィールド名を推測することができます。もう一つの有用な情報は、ファイル名です。そこにどのようなタイプのコードがあるかを推測することができます。
また、ファイル名から既知のオープンソースライブラリを認識することも可能です。

\lstinputlisting[caption=情報豊富なassert()呼び出しの例,style=customasmx86]{digging_into_code/assert_examples.lst}

条件とファイル名の両方を「グーグル」することをお勧めします。これにより、オープンソースライブラリにたどり着くことができます。
例えば、「sp->lzw\_nbits <= BITS\_MAX」を「グーグル」すると、予想通りLZW圧縮に関連するオープンソースコードが見つかります。